\documentclass[11pt]{article}
\newcommand{\routelabel}{Setup and practice}
\usepackage{eastbridge-handout}
\hypersetup{pdftitle={RPS Workshop: Setup and practice}}
\begin{document}
\thispagestyle{plain}

\workshoptitle{A bot in the arena}{Getting started}

\begin{context}
You'll write a Python bot that plays rock, paper, scissors against everyone else's bots. Both players choose at the same time. We'll start by entering the supplied random bot, then spend most of the session changing how it plays.

\end{context}

\subsection*{Start with a working project}

\textbf{Lab machine:} the kit is already in your home folder. Run the commands below. \textbf{Personal laptop:} first follow page 2 (macOS/Linux) or page 3 (Windows), then return here to submit. All levels use this setup handout.

\begin{lstlisting}[style=shell]
cd ~/rps-event-kit
sh ./start.sh ~/rps-bot
cd ~/rps-bot
source .venv/bin/activate
rps-cli test
rps-cli validate
\end{lstlisting}

These commands install the kit and its dependencies from the supplied folder. \textbf{uv and Python 3.11+} are already installed on the lab machines. Windows users should use \code{.\textbackslash{}.venv\textbackslash{}Scripts\textbackslash{}rps-cli.exe} wherever this handout says \code{rps-cli}. Whenever you open a new terminal, return to \code{\textasciitilde{}/rps-bot} and activate \code{.venv} again.


Open \textbf{bot.py} in your editor. It currently chooses a random move, and its tests should pass as supplied. If you accidentally delete a starter file, \code{rps-cli init} will restore missing files. Adding \code{--force} also overwrites files you've edited.


\subsection*{Submit your bot}

Your facilitator will give you the \textbf{server URL, league slug and submit token}. Substitute those values in the commands below, including replacing \code{TOKEN} with the supplied token.


\begin{lstlisting}[style=shell]
rps-cli config set api_url https://arena.eastbrid.ge
rps-cli config set league rps
rps-cli config set token TOKEN
rps-cli doctor
rps-cli submit "Your Team Name"
rps-cli status "Your Team Name"
\end{lstlisting}

The settings are saved in your project folder. Reuse the \textbf{same team name} for later submissions so they replace your earlier version. The server checks each upload before letting it play; \code{status} shows when it becomes active, or why it was rejected. Run \code{submit} again whenever you want the tournament to use your latest edits.


\begin{checkpoint}
\textbf{Before moving on:} run \code{status} and find your active version. If you're stuck on installation or submission, ask a facilitator to take a look.
\end{checkpoint}
\clearpage
\workshoptitle{Bring your own laptop}{macOS and Linux}

\begin{context}
Download the event kit, install uv and Python, and create a folder for your bot.
The ZIP includes the client kit, its dependencies and all handouts. You do not
need Git, a compiler or a copy of the Arena server.
\end{context}

\lessonstep{Step 1}{Download and extract the kit}
Open the \href{https://github.com/Eastbridge-Academy/rps-client-kit/releases/latest}{GitHub release page}
(address below) and choose \textbf{rps-event-kit.zip} under Assets. Extract it
and find the folder containing \textbf{start.sh}, \textbf{wheelhouse} and
\textbf{requirements.txt}. Keep these files together; your bot will live in a
separate folder.

{\small\url{https://github.com/Eastbridge-Academy/rps-client-kit/releases/latest}\par}

\lessonstep{Step 2}{Install uv and Python}
Open Terminal. If \code{uv --version} already works, skip the first command.
Otherwise, install uv:

\begin{lstlisting}[style=shell]
curl -LsSf https://astral.sh/uv/install.sh | sh
\end{lstlisting}

Close Terminal and open it again, then install Python for the workshop:

\begin{lstlisting}[style=shell]
uv python install 3.12
\end{lstlisting}

These downloads need internet access. After this step, installing the kit and
playing local matches use the downloaded files. Submitting your bot and viewing
the tournament still need a connection.

\lessonstep{Step 3}{Create your bot project}
In Terminal, type \code{cd} followed by a space, drag the extracted folder
containing \textbf{start.sh} into the window, and press Enter. On Linux, you can
also right-click the folder and choose Open in Terminal. Then run:

\begin{lstlisting}[style=shell]
sh ./start.sh "$HOME/rps-bot"
cd "$HOME/rps-bot"
source .venv/bin/activate
rps-cli test
rps-cli validate
\end{lstlisting}

Open \textbf{bot.py} in your editor. Select the Python interpreter in
\textbf{rps-bot/.venv} if your editor asks. Each new terminal needs the
\code{cd} and \code{source} commands again.

\begin{checkpoint}
\textbf{Join the event:} return to \emph{Submit your bot} on page 1.
The facilitator supplies the event token. Run all bot commands from your
\textbf{rps-bot} folder.
\end{checkpoint}

\clearpage
\workshoptitle{Bring your own laptop}{Windows 10 / 11}

\begin{context}
Download \textbf{rps-event-kit.zip} from the release page below. Right-click the
ZIP and choose Extract All. Find the folder containing \textbf{start.sh},
\textbf{wheelhouse} and \textbf{requirements.txt}; it may be one folder deeper
than the first folder Windows opens.
\end{context}

{\small\url{https://github.com/Eastbridge-Academy/rps-client-kit/releases/latest}\par}

\lessonstep{Step 1}{Install uv and Python}
Open PowerShell from the Start menu. If \code{uv --version} already works,
skip the installer command. Otherwise, run:

\begin{lstlisting}[style=plaincode]
winget install --id astral-sh.uv -e
\end{lstlisting}

Close PowerShell and open it again, then run:

\begin{lstlisting}[style=plaincode]
uv python install 3.12
\end{lstlisting}

If winget is unavailable, ask a facilitator or use the Windows installer at
\url{https://docs.astral.sh/uv/getting-started/installation/}.
These downloads need internet access; installing the kit and practicing use
local files.

\lessonstep{Step 2}{Create your bot project}
In File Explorer, open the extracted folder containing \textbf{requirements.txt}.
Click the address bar, type \code{powershell}, and press Enter. Run the commands
below; each final backtick continues the command onto the next line.

\begin{lstlisting}[style=plaincode]
$RpsProject = "$HOME\rps-bot"
uv venv --offline --python 3.12 "$RpsProject\.venv"
uv pip install --offline --no-index `
  --python "$RpsProject\.venv\Scripts\python.exe" `
  --find-links .\wheelhouse -r .\requirements.txt
cd $RpsProject
.\.venv\Scripts\rps-cli.exe init
.\.venv\Scripts\rps-cli.exe test
.\.venv\Scripts\rps-cli.exe validate
\end{lstlisting}

Open \textbf{bot.py} in the new \textbf{rps-bot} folder. In your editor, select
\textbf{.venv/Scripts/python.exe} as the Python interpreter.

\lessonstep{Step 3}{Use the handouts}
Where a handout says \code{rps-cli}, use
\code{.\textbackslash{}.venv\textbackslash{}Scripts\textbackslash{}rps-cli.exe}.
Each new PowerShell window needs \code{cd "\$HOME\textbackslash{}rps-bot"}
first. You do not need to activate a PowerShell script or change its execution policy.

\begin{checkpoint}
\textbf{Join the event:} follow \emph{Submit your bot} on page 1, using the
Windows executable path. Reuse the same team name for every upload.
\end{checkpoint}

\clearpage
\workshoptitle{Practice matches}{Scores and local testing}

\begin{context}
You can play practice matches on your own machine while your submitted bot competes in the arena. Save a copy before trying a new strategy, so you can compare the two versions or go back to the earlier one.

\end{context}

\subsection*{Reading the score}

A \textbf{throw} is one simultaneous choice. A \textbf{match} is a fixed series, normally 501 throws. A \textbf{league round} schedules each pair of bots. Despite the historical name \code{best\_of}, the arena plays the full series; it does not stop at 251 wins. Drawn throws still count, and a 501-throw match can finish level.


For example, a 501-throw match ending 290-102 had 109 draws. The rating system records a match win for the first bot, regardless of the size of that margin. Local practice reports match results and throw totals; ratings belong to the arena.


\begin{lstlisting}[style=shell,numbers=none,xleftmargin=4pt,framexleftmargin=0pt]
rps-cli play --against random_uniform --games 5 --seed 100
\end{lstlisting}

\begin{checkpoint}
\textbf{Comparing versions:} use the same opponents, seeds and series length for both bots. Once you've chosen a version, try it on some fresh seeds as well. Check the error counts before you read much into the score.
\end{checkpoint}

\subsection*{What gets uploaded}

Your upload includes the installed \code{rpsdk}, imported local helper modules and packages, and files in \code{data/}. Other packages from your virtual environment are left out, so use the standard library for these exercises. Open data files relative to \code{Path(\_\_file\_\_).parent}. Each move has a time limit; network calls, sleeps and large training jobs are likely to exceed it.


\subsection*{During the session}

Allow about 20 minutes for setup and the rules, then 70 minutes for your chosen route. Leave another 20 minutes to test your final version on fresh seeds and 10 to submit it. You can skip an optional exercise or switch routes whenever you like.
\par

\end{document}
